% id: 1093
% book: RPM 대수 · 10 수학적 귀납법
% section: 유형 08 수학적 귀납법
% figure: none
% answer: ④
% answer_source: 답지
% variant_level: 0 (원본 전사)
% 이 파일은 build.mjs 가 items.json 에서 만든다. 고칠 때는 items.json 을 고친다.
\hspace*{2mm}\begin{minipage}[t]{\dimexpr\linewidth-2mm\relax}\vspace{0pt}\smash{\raisebox{1.5mm}{\dmkichul{RPM 대수 1093번 · 149쪽}}}
모든 자연수 $n$에 대하여 명제 $p(n)$이 참이면 명제 $p(n+2)$가 참일 때, \textbf{보기}에서 항상 옳은 것만을 있는 대로 고른 것은? \cond{$k$는 자연수이다.} \bogi{ㄱ. $p(1)$이 참이면 $p(2k+1)$이 참이다.\\ ㄴ. $p(2)$가 참이면 $p(2k+3)$이 참이다.\\ ㄷ. $p(1)$, $p(2)$가 참이면 $p(k)$가 참이다.}
\begin{choices32}{ㄱ}{ㄷ}{ㄱ, ㄴ}{ㄱ, ㄷ}{ㄱ, ㄴ, ㄷ}\end{choices32}\end{minipage}
